\pagestyle{empty}
\tight
\begin{tblr}{width=\textwidth,colspec={|Q[c,m,1.9cm]|X[l,m]|},hlines,vlines,hspan=minimal,row{1}={bg=headblue},rowsep=1pt,colsep=3pt}
\SetCell{c,m}\hfil\logo\hfil & \textbf{POLITEKNIK NEGERI MALANG}\par\textbf{JURUSAN TEKNOLOGI INFORMASI}\par\textbf{PROGRAM STUDI: D4 TEKNIK INFORMATIKA} \\
\SetCell[c=2]{c} \textbf{RENCANA PEMBELAJARAN SEMESTER (RPS)} & \\
\end{tblr}
\begin{longtblr}{width=\textwidth,colspec={|Q[l,m,1.9cm]|X[0.85,l,m]|X[1.45,l,m]|X[0.75,l,m]|X[1.15,l,m]|},hlines,vlines,hspan=minimal,rowsep=1pt,colsep=2pt,row{1,3}={bg=headgray,font=\bfseries}}
MATA KULIAH & KODE & BOBOT (SKS)/jam & SEMESTER & TGL. PENYUSUNAN \\
Algoritma dan Struktur Data & RTI252008 & 2 SKS/4 jam & 2 & 31 Januari 2026 \\
OTORISASI & Dosen Pengembang RPS & Koordinator RMK & \SetCell[c=2]{l} Ka PRODI &  \\
 & Mungki Astiningrum, S.T., M.Kom.\par Mamluatul Hani'ah, S.Kom., M.Kom.\par Vivi Nur Wijayaningrum, S.Kom., M.Kom. &  & \SetCell[c=2]{l}{\vspace{8mm}\par Dr. Ely Setyo Astuti, S.T., M.T.} &  \\
\SetCell[r=9]{l} Capaian Pembelajaran (CP) & \SetCell[c=4]{l,bg=headgray,font=\bfseries} Capaian Pembelajaran Lulusan Program Studi (CPL-Prodi) &  &  &  \\
 & CPL02 & \SetCell[c=3]{l} Menguasai konsep teoritis bidang pengetahuan mengenai berbagai prinsip rekayasa sistem/perangkat lunak dalam merancang solusi aplikasi multi-platform yang relevan dengan kebutuhan stakeholder. &  &  \\
 & CPL09 & \SetCell[c=3]{l} Memiliki pemahaman yang memadai terkait cara kerja sistem komputer dan berbagai teknik/pendekatan berbasis komputasi untuk memecahkan masalah stakeholder. &  &  \\
 & \SetCell[c=4]{l,bg=headgray,font=\bfseries} Capaian Pembelajaran mata kuliah (CPL-MK) &  &  &  \\
 & CPMK0207 & \SetCell[c=3]{l} Mampu menganalisis kompleksitas permasalahan komputasi kemudian memilih dan merancang struktur data serta algoritma yang optimal. &  &  \\
 & CPMK0903 & \SetCell[c=3]{l} Mahasiswa mampu menjelaskan prinsip logika komputasi serta menerapkan teknik dasar pemrograman dan algoritma untuk menganalisis dan menyelesaikan permasalahan. &  &  \\
 & \SetCell[c=4]{l,bg=headgray,font=\bfseries} Kemampuan akhir tiap tahapan belajar (Sub-CPMK) &  &  &  \\
 & SCPMK0207-01701 & \SetCell[c=3]{l} Mahasiswa mampu menjelaskan karakteristik berbagai struktur data serta alur pengelolaan data dalam mendukung penyelesaian permasalahan komputasi. &  &  \\
 & SCPMK0903-01701 & \SetCell[c=3]{l} Mahasiswa mampu menganalisis permasalahan komputasi pada struktur data berbasis array of object, serta menjelaskan dan membandingkan pendekatan algoritmik yang digunakan untuk pengolahan dan pencarian data. &  &  \\
\SetCell[r=2]{l} Penilaian & \SetCell[c=4]{l} *Nilai CPMK didapat dari rata-rata nilai SUB-CPMK yang dibebankan &  &  &  \\
 & \SetCell[c=4]{l}{\tiny\begin{tblr}{width=\linewidth,colspec={|Q[c,m,0.08\linewidth]|Q[c,m,0.14\linewidth]|X[l,m]|X[l,m]|X[l,m]|X[l,m]|X[l,m]|X[l,m]|Q[c,m,0.07\linewidth]|},hlines,vlines,hspan=minimal,rowsep=1pt,colsep=0.7pt,row{1,2}={bg=headgray,font=\bfseries}}
Kode CPMK & Kode SCPMK & \SetCell[c=6]{c} Bobot per Bentuk Penilaian & & & & & & Total Bobot \\
 & & Tugas & Tes tulis 1\par Asessment 1\par Kuis 1 & Tes tulis 2\par Asessment 2\par UTS & Tes tulis 3\par Asessment 1\par Kuis 2 & Tes tulis 4\par Asessment 2\par UAS & PBL & \\
CPMK 207 & SCPMK0207-01701 & 10\% Materi\par 9,10,11,12,14,15 & & & 20\% Materi\par 9,10,11,12 & 20\% Materi\par 14,15 & & 50\% \\
CPMK 903 & SCPMK0903-01701 & 10\% Materi\par 1,2,3,5,6,7 & 20\% Materi\par 1,2,3 & 20\% Materi\par 5,6,7 & & & & 50\% \\
\SetCell[c=8]{r}\textbf{Total Per Penilaian} & & & & & & & & \textbf{100\%} \\
\end{tblr}} &  &  &  \\
Diskripsi Singkat Mata Kuliah & \SetCell[c=4]{l} Matakuliah Algoritma dan Struktur Data memberikan pengetahuan dan keterampilan pembuatan algoritma dan struktur data yang meliputi brute force, divide-conquer, stack, queue, linked list, tree, graph, sorting, dan searching. &  &  &  \\
Bahan Kajian / Materi Pembelajaran & \SetCell[c=4]{l}\begin{enumerate}\item Searching\item Sorting\item Queue\item Stack\item Linked List\item Tree\item Graph\item Brute Force\item Divide-Conquer\item Depth First Search\item Breadth First Search\end{enumerate} &  &  &  \\
\SetCell[r=2]{l} Pustaka & \SetCell[c=4]{l} \textbf{Utama:}\begin{enumerate}\item Goodrich, M. T., Tamassia, R., and Goldwasser, M. H. 2014. Data Structures \& Algorithms in Java, 6th Edition. Wiley Global Education.\item Ramadhani, C. 2015. Dasar Algoritma dan Struktur Data dengan Bahasa Java. Yogyakarta: Andi Publisher.\item Nugroho, A. 2008. Algoritma dan Struktur Data dalam Bahasa Java. Yogyakarta: Andi Publisher.\item Hariyanto, B. 2007. Struktur Data. Bandung: Informatika.\end{enumerate} &  &  &  \\
 & \SetCell[c=4]{l} \textbf{Pendukung:}\begin{enumerate}\item Munir, R. 2015. Algoritma dan Pemrograman. Bandung: Informatika.\end{enumerate} &  &  &  \\
Media Pembelajaran & \SetCell[c=2]{l} \textbf{Software:} IDE Java, compiler/JDK, LMS, dan perangkat presentasi. &  & \SetCell[c=2]{l} \textbf{Hardware:} Komputer, laptop, dan proyektor. &  \\
Nama Dosen Pengampu & \SetCell[c=4]{l}\begin{enumerate}\item Mamluatul Hani'ah, S.Kom., M.Kom.\item Vivi Nur Wijayaningrum, S.Kom., M.Kom.\item Imam Fahrur Rozi, S.T., M.T.\item Novta Dany'el Irawan, S.Pd., M.T.\item Mungki Astiningrum, S.T., M.Kom.\item Irsyad Arif Mashudi, S.Kom., M.Kom.\item Rokhimatul Wakhidah, S.Pd., M.T.\end{enumerate} &  &  &  \\
Matakuliah Syarat & \SetCell[c=4]{l} - &  &  &  \\
\end{longtblr}
\begin{longtblr}{width=\textwidth,colspec={|Q[c,m,0.06\textwidth]|X[1.35,l,m]|X[1.08,l,m]|X[1.16,l,m]|X[0.7,l,m]|X[1.24,l,m]|X[1.06,l,m]|X[0.98,l,m]|Q[c,m,0.05\textwidth]|},hlines,vlines,hspan=minimal,rowhead=2,row{1,2}={bg=headgreen,font=\bfseries},rowsep=1pt,colsep=1.5pt}
Mgg.\par Ke & Kemampuan Akhir Yang Direncanakan (Sub-CP-MK) & Bahan kajian (Materi Pembelajaran) & Bentuk dan Metode Pembelajaran & Estimasi Waktu & Pengalaman Belajar Mahasiswa & Kriteria \& Bentuk Penilaian & Indikator Penilaian & Bobot Penilaian (\%) \\
(1) & (2) & (3) & (4) & (5) & (6) & (7) & (8) & (9) \\
1 & SCPMK0903-01701 & Pemilihan, perulangan, array, dan fungsi. & Modalitas: Blended Learning. Bentuk: Luring. Metode: diskusi kelompok dan studi kasus. Sumber belajar: LMS Polinema. & 2 x 2 x 100' tatap muka dan tugas terstruktur; 2 x 2 x 70' tugas mandiri; penugasan 1 x 50'. & Mahasiswa memahami konsep dasar algoritma, mengimplementasikan flowchart untuk studi kasus, dan mempresentasikan hasil. & Rubrik kriteria penilaian; keaktifan diskusi; tes tulis atau tugas pemecahan studi kasus dengan flowchart. & Pemahaman konsep dasar algoritma; ketepatan pembuatan flowchart untuk menyelesaikan studi kasus. & 1.7 \\
2 & SCPMK0903-01701 & Konsep objek, class, atribut, method, instansiasi, konstruktor, dan class diagram. & Modalitas: Blended Learning. Bentuk: Luring. Metode: diskusi kelompok dan studi kasus. Sumber belajar: LMS Polinema. & 2 x 2 x 100' tatap muka dan tugas terstruktur; 2 x 2 x 70' tugas mandiri; penugasan 1 x 50'. & Mahasiswa memahami object, mendeklarasikan class, atribut, method, membuat object, mengakses atribut dan method, serta menuangkan object ke class diagram. & Rubrik kriteria penilaian; keaktifan diskusi; tes tulis atau tugas pemecahan studi kasus. & Pemahaman konsep object, class, atribut, dan method; ketepatan pembuatan class diagram. & 1.7 \\
3 & SCPMK0903-01701 & Array of object: deklarasi, instansiasi, assignment, dan penampilan data. & Modalitas: Blended Learning. Bentuk: Kuliah luring. Metode: diskusi kelompok virtual dan studi kasus. Sumber belajar: LMS Polinema. & 2 x 2 x 100' tatap muka dan tugas terstruktur; 2 x 2 x 70' tugas mandiri; penugasan 1 x 50'. & Mahasiswa memahami konsep array of object dan menerapkan array object dalam flowchart untuk berbagai studi kasus. & Rubrik kriteria penilaian; keaktifan diskusi; tes tulis atau tugas pemecahan studi kasus dengan flowchart. & Pemahaman array of object; ketepatan flowchart sebagai implementasi algoritma untuk menyelesaikan studi kasus. & 1.7 \\
4 & SCPMK0903-01701 & Materi pekan 1-3. & Ujian teori pilihan ganda dan esai, closed book. & Sesuai jadwal asesmen. & Mahasiswa mengerjakan soal ujian materi pekan 1-3 secara mandiri. & Ujian. & Ketepatan jawaban dengan kunci jawaban. & 20 \\
5 & SCPMK0903-01701 & Algoritma brute force, divide-conquer, kompleksitas Big-O, dan perhitungan notasi Big-O. & Modalitas: Blended Learning. Bentuk: Kuliah luring. Metode: diskusi kelompok virtual dan studi kasus. Sumber belajar: LMS Polinema. & 2 x 2 x 100' tatap muka dan tugas terstruktur; 2 x 2 x 70' tugas mandiri; penugasan 1 x 50'. & Mahasiswa memahami brute force dan divide-conquer serta mengimplementasikannya dalam flowchart. & Rubrik kriteria penilaian; keaktifan diskusi; tes tulis atau tugas pemecahan studi kasus dengan flowchart. & Pemahaman brute force dan divide-conquer; ketepatan flowchart sebagai implementasi algoritma. & 1.7 \\
6 & SCPMK0903-01701 & Sorting: bubble sort, selection sort, dan insertion sort. & Modalitas: Blended Learning. Bentuk: Kuliah luring. Metode: diskusi kelompok virtual dan studi kasus. Sumber belajar: LMS Polinema. & 2 x 2 x 100' tatap muka dan tugas terstruktur; 2 x 2 x 70' tugas mandiri; penugasan 1 x 50'. & Mahasiswa memahami konsep algoritma sorting, menentukan algoritma yang digunakan, dan membuat flowchart untuk studi kasus sorting. & Rubrik kriteria penilaian; keaktifan diskusi; tes tulis atau tugas pemecahan studi kasus dengan flowchart. & Pemahaman algoritma sorting; ketepatan flowchart sebagai implementasi algoritma untuk menyelesaikan studi kasus. & 1.7 \\
7 & SCPMK0903-01701 & Sequential search, binary search, dan pengayaan merge sort. & Modalitas: Blended Learning. Bentuk: Luring. Metode: diskusi kelompok dan studi kasus. Sumber belajar: LMS Polinema. & 2 x 2 x 100' tatap muka dan tugas terstruktur; 2 x 2 x 70' tugas mandiri; penugasan 1 x 50'. & Mahasiswa memahami konsep searching, sequential search, binary search, dan melakukan simulasi manual searching. & Ketepatan dan penguasaan; keaktifan diskusi; ketepatan jawaban tugas minggu 7. & Pemahaman penerapan searching; kesesuaian jawaban tugas. & 1.7 \\
8 & SCPMK0903-01701 & Materi pekan 5, 6, dan 7. & Ujian teori pilihan ganda dan esai, closed book. & Sesuai jadwal asesmen. & Mahasiswa mengerjakan soal ujian materi pekan 5-7 secara mandiri. & Ujian. & Ketepatan jawaban dengan kunci jawaban. & 20 \\
9 & SCPMK0207-01701 & Konsep struktur data stack, operasi stack, dan postfix expressions. & Modalitas: Blended Learning. Bentuk: Luring. Metode: Contextual Teaching and Learning (CTL). Sumber belajar: LMS Polinema. & 2 x 2 x 100' tatap muka dan tugas terstruktur; 2 x 2 x 70' tugas mandiri; penugasan 1 x 50'. & Mahasiswa memahami struktur data stack, operasi pada stack, dan mengimplementasikan konsep stack untuk menyelesaikan studi kasus. & Ketepatan dan penguasaan; keaktifan diskusi; ketepatan jawaban tugas minggu 9. & Pemahaman penerapan konsep stack; kesesuaian jawaban tugas. & 1.7 \\
10 & SCPMK0207-01701 & Konsep struktur data queue dan operasi-operasi pada queue. & Modalitas: Blended Learning. Bentuk: Luring. Metode: Contextual Teaching and Learning (CTL). Sumber belajar: LMS Polinema. & 2 x 2 x 100' tatap muka dan tugas terstruktur; 2 x 2 x 70' tugas mandiri; penugasan 1 x 50'. & Mahasiswa memahami struktur data queue, operasi pada queue, dan mengimplementasikan konsep queue untuk menyelesaikan studi kasus. & Ketepatan dan penguasaan; keaktifan diskusi; ketepatan jawaban tugas minggu 10. & Pemahaman penerapan konsep queue; kesesuaian jawaban tugas. & 1.7 \\
11 & SCPMK0207-01701 & Linked list. & Modalitas: Blended Learning. Bentuk: Kuliah. Metode: diskusi kelompok dan studi kasus. Sumber belajar: LMS Polinema. & 2 x 2 x 100' tatap muka dan tugas terstruktur; 2 x 2 x 70' tugas mandiri; penugasan 1 x 50'. & Mahasiswa menerapkan konsep linked list dalam diagram flowchart dan menunjukkan hasil pemahaman konsep linked list. & Rubrik kriteria penilaian; presentasi; tes tulis atau tugas pemecahan studi kasus dengan flowchart. & Pemahaman konsep linked list; ketepatan diagram flowchart sesuai studi kasus; performa presentasi. & 1.7 \\
12 & SCPMK0207-01701 & Double linked list. & Modalitas: Blended Learning. Bentuk: Kuliah. Metode: diskusi kelompok dan studi kasus. Sumber belajar: LMS Polinema. & 2 x 2 x 100' tatap muka dan tugas terstruktur; 2 x 2 x 70' tugas mandiri; penugasan 1 x 50'. & Mahasiswa memahami double linked list sebagai pengembangan linked list, menerapkannya dalam diagram flowchart, dan menunjukkan hasil pemahaman. & Rubrik kriteria penilaian; presentasi; tes tulis atau tugas pemecahan studi kasus dengan flowchart. & Pemahaman konsep double linked list; ketepatan diagram flowchart sesuai studi kasus; performa presentasi. & 1.7 \\
13 & SCPMK0207-01701 & Materi pekan 9-12. & Ujian teori pilihan ganda dan/atau esai, closed book. & Sesuai jadwal asesmen. & Mahasiswa mengerjakan soal ujian materi pekan 9-12 secara mandiri. & Ujian. & Ketepatan jawaban dengan kunci jawaban. & 20 \\
14 & SCPMK0207-01701 & Tree dan binary search tree. & Modalitas: Blended Learning. Bentuk: Kuliah. Metode: diskusi kelompok dan studi kasus. Sumber belajar: LMS Polinema. & 2 x 2 x 100' tatap muka dan tugas terstruktur; 2 x 2 x 70' tugas mandiri; penugasan 1 x 50'. & Mahasiswa memahami konsep tree, penerapan binary tree, konsep binary search tree, dan tahapan penerapannya. & Ketepatan dan penguasaan; keaktifan diskusi; ketepatan jawaban tugas minggu 14. & Pemahaman konsep tree dan binary search tree; ketepatan menjawab studi kasus. & 1.7 \\
15 & SCPMK0207-01701 & Graph. & Modalitas: Blended Learning. Bentuk: Kuliah. Metode: diskusi kelompok dan studi kasus. Sumber belajar: LMS Polinema. & 2 x 2 x 100' tatap muka dan tugas terstruktur; 2 x 2 x 70' tugas mandiri; penugasan 1 x 50'. & Mahasiswa membuat algoritma graph secara umum dan menerapkan algoritma graph pada program. & Ketepatan dan penguasaan; keaktifan diskusi; ketepatan jawaban tugas minggu 15. & Pemahaman konsep graph; ketepatan menjawab tugas. & 1.7 \\
16 & SCPMK0207-01701 & Materi pekan 14 dan 15. & Ujian teori pilihan ganda dan/atau esai, closed book. & Sesuai jadwal asesmen. & Mahasiswa mengerjakan soal ujian materi pekan 14-15 secara mandiri. & Ujian. & Ketepatan jawaban dengan kunci jawaban. & 20 \\
17 & SCPMK0207-01701, SCPMK0903-01701 & Materi pekan 1-15. & Pengayaan untuk mahasiswa sesuai Sub-CPMK yang belum tercapai. & Sesuai kebutuhan. & Mahasiswa mengerjakan pengayaan materi pekan 1-15 secara mandiri. & Ujian atau pengayaan. & Ketepatan jawaban dengan kunci jawaban. & 0 \\
\end{longtblr}
